\chapter{Model Map Reduce}
% ─────────────────────────────────────────────────────────────────────
\section{Motivazione e Contesto}
% ─────────────────────────────────────────────────────────────────────

\subsection{Ingredienti fondamentali dei sistemi Big Data}

La maggior parte dei framework per l'elaborazione di Big Data si basa su
\textbf{tre ingredienti principali}:

\begin{enumerate}
  \item \textbf{Infrastruttura di storage distribuito} -- I \emph{Distributed
    File System} (DFS) garantiscono storage scalabile e fault-tolerant.
  \item \textbf{Modello computazionale / di programmazione} -- MapReduce
    come modello per esprimere computazioni data-parallel su larga scala,
    progettato per girare su cluster, tipicamente su dati residenti nel DFS.
  \item \textbf{Framework di esecuzione / orchestrazione} -- Responsabile
    di scheduling, sincronizzazione, data movement (shuffle), fault tolerance
    e monitoring.
\end{enumerate}

\begin{takeaway}
  DFS + modello data-parallel (MapReduce) + runtime fanno sì che un datacenter
  si comporti come una \emph{``singola macchina facile da usare''} per i Big Data.
\end{takeaway}

% ─────────────────────────────────────────────────────────────────────
\subsection{Tecniche di decomposizione per il calcolo parallelo}

MapReduce è un modello per descrivere come avviene il calcolo nei cluster.
Lo usiamo per \textbf{progettare e analizzare algoritmi paralleli} con garanzie
che funzionino bene in pratica.

Il modello si ispira alle classiche tecniche di decomposizione che suddividono
un problema in sottoproblemi risolvibili in parallelo.

\begin{importante}
  \textbf{Idea chiave:} più i sottoproblemi sono indipendenti, maggiore è
  il potenziale speedup se elaborati in parallelo.
\end{importante}

\paragraph{Divide and Conquer}

La strategia classica di decomposizione:
\begin{itemize}
  \item \textbf{Divide:} suddivide il lavoro in molti sottoproblemi.
  \item \textbf{Conquer:} risolve i sottoproblemi su worker/unità computazionali
    diverse.
  \item \textbf{Combine:} unisce i risultati parziali nell'output finale.
\end{itemize}

\paragraph{Input-Data Decomposition}

\begin{itemize}
  \item Partizione dei \textbf{dati di input} tra i worker.
  \item Ogni worker elabora la propria porzione locale \textbf{indipendentemente}.
  \item Esempio: contare le parole in una collezione di documenti -- ogni
    worker processa un sottoinsieme di documenti.
  \item La \textbf{fase Map} di MapReduce è un esempio canonico di
    input-data decomposition.
\end{itemize}

\paragraph{Output-Data Decomposition}

\begin{itemize}
  \item Partizione dello \textbf{spazio di output} tra i worker.
  \item Ogni worker è responsabile della produzione di una porzione dell'output.
  \item Esempio: calcolo di un istogramma -- ogni worker è responsabile di
    un sottoinsieme di bin.
  \item La \textbf{fase Reduce} di MapReduce usa output-data decomposition
    (ogni reducer produce output per le chiavi assegnategli).
\end{itemize}

\paragraph{Owner-Computes-Rule Decomposition}

\begin{itemize}
  \item Il worker che \emph{possiede} un dato è responsabile del calcolo su di esso.
  \item Minimizza il movimento di dati \textbf{co-localizzando} computazione e dati.
  \item Esempio: operazioni matriciali in cui ogni worker possiede un blocco
    della matrice.
\end{itemize}

\paragraph{Exploratory Decomposition}

\begin{itemize}
  \item Utile per problemi di ricerca in cui lo spazio delle soluzioni è esplorato
    in parallelo.
  \item Worker diversi esplorano rami diversi dell'albero di ricerca.
  \item Esempio: branch-and-bound parallelo, ricerca su albero di gioco.
\end{itemize}

\paragraph{Speculative Decomposition}

\begin{itemize}
  \item Esegue più computazioni possibili in parallelo \emph{prima} di sapere
    quale è necessaria.
  \item Scarta i risultati che risultano inutili.
  \item Esempio: \textbf{backup task} in MapReduce (esecuzione speculativa
    dei task lenti).
\end{itemize}

\paragraph{MapReduce combina più strategie}

\begin{tabular}{@{}ll@{}}
  \toprule
  \textbf{Elemento MapReduce} & \textbf{Strategia} \\
  \midrule
  Fase Map & Input-data decomposition \\
  Fase Reduce & Output-data decomposition \\
  Data locality & Owner-computes-rule \\
  Backup task & Speculative decomposition \\
  \bottomrule
\end{tabular}

\vspace{0.5em}
\begin{importante}
  La comunicazione di rete è spesso il collo di bottiglia: si mira a
  \textbf{minimizzare il movimento dei dati} e sfruttare la \textbf{data locality}.
\end{importante}

% ─────────────────────────────────────────────────────────────────────
\subsection{Efficacia della decomposizione}

L'efficacia dipende da:
\begin{itemize}
  \item \textbf{Struttura del problema:} i sottoproblemi sono risolvibili
    indipendentemente?
  \item \textbf{Load balance:} i worker ricevono quantità simili di lavoro?
  \item \textbf{Costo di comunicazione:} quanti dati devono essere trasmessi
    sulla rete (shuffle)?
  \item \textbf{Effetti di sistema:} fallimenti e \emph{stragglers} (task lenti
    che ritardano il completamento).
\end{itemize}

In pratica riutilizziamo idee degli algoritmi paralleli classici, ma a scala
dove comunicazione e I/O possono dominare sul calcolo puro.

% ─────────────────────────────────────────────────────────────────────
\subsection{Esempi di problemi a scala Web}

\begin{itemize}
  \item \textbf{Web page ranking (PageRank):} moltiplicazioni matrice-vettore
    ripetute su matrici sparse enormi; ogni worker elabora un sottoinsieme
    di righe/archi.
  \item \textbf{Mining da social network:} calcolo di proprietà su grafi con
    miliardi di nodi/archi; computazioni locali su sottografi + aggregazione
    di statistiche globali.
  \item \textbf{Rilevamento di plagi / near-duplicate:} l'approccio naive
    confronta tutte le coppie di documenti (infeasible); metodi scalabili
    generano coppie candidate e le verificano.
  \item \textbf{Workload di training:} conteggio di frequenze, co-occorrenze
    o gradienti su dataset enormi; risultati parziali locali + riduzioni globali
    (somma, media, max, \ldots).
\end{itemize}

% ─────────────────────────────────────────────────────────────────────
\subsection{Come i Massively Parallel Computational Models (MPCM) differiscono
  dai modelli classici}

Rispetto ai modelli paralleli classici, gli MPCM assumono tipicamente:
\begin{itemize}
  \item \textbf{Shared-nothing:} nessuna memoria condivisa tra i worker.
  \item \textbf{Dati partizionati} (ed eventualmente replicati) tra le macchine.
  \item \textbf{Controllo limitato} su placement dello storage e risorse
    computazionali: gli algoritmi devono funzionare senza conoscere il
    placement dei dati (\emph{agnostic} alla distribuzione) e senza conoscere
    il numero di unità computazionali.
\end{itemize}

Queste assunzioni migliorano \textbf{portabilità} e \textbf{robustezza}:
\begin{itemize}
  \item Preferenza per esecuzione \emph{data-local}.
  \item Lo stesso programma gira su cluster di dimensioni/topologie diverse.
  \item Scheduling e fault tolerance gestiti dal sistema.
\end{itemize}

% ─────────────────────────────────────────────────────────────────────
\subsection{Perché MapReduce è importante}

MapReduce (e i suoi discendenti) ha influenzato molti sistemi moderni:
\begin{itemize}
  \item Implementazioni interne di Google.
  \item \textbf{Apache Hadoop} (MapReduce open-source + ecosistema).
  \item Apache Hive e altri strumenti SQL-on-Hadoop.
  \item Servizi cloud gestiti (Google Cloud Dataproc, AWS EMR, Azure HDInsight).
  \item \textbf{Apache Spark} (va oltre MapReduce, ma ne eredita le idee chiave).
\end{itemize}

\begin{takeaway}
  Il motivo chiave del successo: il modello è abbastanza semplice da supportare
  astrazioni di alto livello e blocchi riutilizzabili, non solo codice distribuito
  di basso livello.
\end{takeaway}

Il successo di MapReduce deriva dalla sua \textbf{semplicità e rigidità}:
il modello computazionale è fisso, quindi relativamente facile da usare.
La progettazione di algoritmi è semplificata quando le assunzioni di modellazione
sono fisse e realistiche. Solo un piccolo insieme di componenti è personalizzabile;
il programmatore prepara i dati e implementa le funzioni.

% ─────────────────────────────────────────────────────────────────────
\newpage
\section{Il Modello MapReduce}
% ─────────────────────────────────────────────────────────────────────

\subsection{Principio ispiratore: spostare il calcolo verso i dati}

\begin{importante}
  \textbf{Motto:} trattare il datacenter come una singola macchina.
  \begin{itemize}
    \item I dati sono memorizzati in modo distribuito dove vengono prodotti.
    \item Replicati per fault tolerance.
    \item Quando è necessaria l'elaborazione, i task vengono schedulati
      \textbf{vicino ai dati} (\emph{data locality}).
  \end{itemize}
  \textbf{Perché:} a scala, la larghezza di banda di rete è preziosa; è spesso
  più veloce \emph{spostare il calcolo verso i dati} che spostare i dati verso
  il calcolo.
\end{importante}

% ─────────────────────────────────────────────────────────────────────
\subsection{Terminologia}

\begin{tabular}{@{}lp{9cm}@{}}
  \toprule
  \textbf{Termine} & \textbf{Significato} \\
  \midrule
  Cluster & Infrastruttura fisica (macchine + rete). \\
  MPCM    & Modello astratto per ragionare su computazione/comunicazione
             in un cluster. \\
  Worker  & Macchina (o processo worker) che esegue task e scambia dati. \\
  Task    & Unità di lavoro schedulata (map task o reduce task). \\
  Chunk / Input Split & Partizione dell'input assegnata a un singolo map task. \\
  Mapper  & Risultato dell'applicazione della funzione map a un singolo
             elemento di input. \\
  Reducer & Risultato dell'applicazione della funzione reduce a una singola
             chiave e ai suoi valori. \\
  \bottomrule
\end{tabular}

% ─────────────────────────────────────────────────────────────────────
\subsection{Modello dei dati: coppie chiave-valore}

\begin{definizione}[-- Coppia chiave-valore]
  Tutta la rappresentazione, l'elaborazione e la memorizzazione in MapReduce
  avviene tramite \textbf{coppie chiave-valore}, denotate:
  \[
    \kv{k}{v}
  \]
  dove $k$ è la chiave e $v$ è il valore. Chiavi e valori possono essere di
  molti tipi: interi, stringhe, tuple, oggetti serializzati, \ldots
\end{definizione}

Questa astrazione è il fondamento per:
\begin{itemize}
  \item Partizionamento e raggruppamento (\emph{shuffle}).
  \item Elaborazione parallela (molti task su molte coppie).
  \item Interfaccia pulita tra codice utente (Map/Reduce) e runtime.
\end{itemize}

\textbf{Perché le chiavi sono importanti:}
\begin{itemize}
  \item Distribuire ed spostare dati nel cluster (hashing): le chiavi
    determinano \emph{dove} i dati vengono memorizzati o trasmessi.
  \item Schedulare i task sui nodi appropriati: le chiavi determinano
    \emph{quale nodo} elabora quali dati.
\end{itemize}

\paragraph{Rappresentazione dei dati di input}

Gli elementi di input elaborati sono coppie chiave-valore:
\begin{itemize}
  \item Un \textbf{elemento} può rappresentare molte cose: una tupla,
    un documento, un record, \ldots
  \item Un \textbf{chunk} (o input split) è una collezione di elementi
    (tipicamente assegnata a un singolo map task).
  \item Se i dati grezzi non sono naturalmente indicizzati, il sistema
    (o l'utente) può assegnare una chiave tecnica (es.\ ID documento,
    byte offset, numero di record).
  \item Quando la chiave è irrilevante, la notazione è $\kv{*}{v}$.
\end{itemize}

\begin{importante}
  \textbf{Vincolo del modello:} nessun elemento di input è condiviso tra due
  chunk (macchine). Ogni elemento è elaborato da esattamente un map task;
  la suddivisione dell'input in chunk è una \textbf{partizione} del dataset.
\end{importante}

% ─────────────────────────────────────────────────────────────────────
\subsection{Le tre fasi di un round MapReduce}

Una computazione MapReduce (``round'') consiste nelle seguenti fasi:

\begin{enumerate}
  \item \textbf{Lettura dell'input} (eventualmente grande):
    ingestione nel DFS se non già distribuito.
  \item \textbf{Fase Map:} converte ogni elemento di input in una o più
    coppie chiave-valore intermedie $\kv{k'}{v'}$.
    Output: un multiset distribuito di coppie $\kv{\mathit{key}}{\mathit{value}}$ intermedie.
  \item \textbf{Fase Shuffle/Group-by-Key:} raggruppa le coppie intermedie
    per chiave, spostando i dati sulla rete affinché tutti i valori con la
    stessa chiave si trovino sulla stessa macchina.
    Output: coppie $\kv{k}{\Lk}$ (chiave -- lista di valori).
  \item \textbf{Fase Reduce:} applica la funzione reduce a ogni coppia
    $\kv{k}{\Lk}$.
    Output: zero o più coppie chiave-valore di output.
  \item \textbf{Scrittura dell'output} (eventualmente grande):
    memorizzazione nel DFS o in un sink esterno.
\end{enumerate}

\begin{nota}
  I passi 1 e 5 sono spesso impliciti: i dati possono essere già nel DFS
  e l'output rimanere distribuito. Questo è l'essenza del
  \emph{``spostare il calcolo verso i dati''}.
\end{nota}

\begin{takeaway}
  La computazione \textbf{alterna} tra lavoro parallelo locale
  (computazione map/reduce) e comunicazione globale (shuffle).
\end{takeaway}

% ─────────────────────────────────────────────────────────────────────
\subsection{Funzione Map: definizione formale}

\begin{definizione}[-- Funzione Map]
  \[
    \forall\, \kv{k}{v} \in I \quad \MAP\!\bigl(\kv{k}{v}\bigr)
      \;\to\; \kv{k'}{v'}^{*}
  \]
  \begin{itemize}
    \item \textbf{Input:} coppie $\kv{k}{v} \in I$ (porzioni diverse di $I$
      sono elaborate da map task diversi).
    \item \textbf{Una chiamata per coppia:} output è una sequenza
      (possibilmente vuota) di coppie intermedie $\kv{k'}{v'}$.
    \item Le chiavi intermedie \textbf{non sono necessariamente univoche}:
      molte coppie possono condividere la stessa chiave.
  \end{itemize}
\end{definizione}

\begin{nota}
  La rappresentazione $\kv{k}{\Lk}$ (raggruppata per chiave) è prodotta
  \emph{dopo} dalla fase shuffle/group, non dalla funzione map stessa.
\end{nota}

La \textbf{Fase Map} è l'esecuzione dei map task su tutti gli input; termina
quando l'ultimo map task termina. Output: unione di tutte le coppie intermedie
prodotte da tutti i map task (distribuita tra i nodi).

% ─────────────────────────────────────────────────────────────────────
\subsection{Funzione Reduce: definizione formale}

\begin{definizione}[-- Funzione Reduce]
  \[
    \forall\, \kv{k}{\Lk} \in I' \quad
    \REDUCE\!\bigl(\kv{k}{\Lk}\bigr) \;\to\; \kv{k''}{v''}^{*}
  \]
  Notazione alternativa:
  \[
    \forall\, \kv{k}{\langle v'\rangle^{*}} \in I' \quad
    \REDUCE\!\bigl(\kv{k}{\langle v'\rangle^{*}}\bigr) \;\to\; \kv{k''}{v''}^{*}
  \]
  \begin{itemize}
    \item \textbf{Input:} coppie raggruppate per chiave $\kv{k}{\Lk} \in I'$;
      $\Lk$ contiene tutti i valori prodotti dalla fase map con chiave $k$.
    \item \textbf{Una chiamata per chiave $k$:} tutti i valori in $\Lk$ sono
      elaborati insieme.
    \item \textbf{Output:} zero o più coppie chiave-valore
      (non necessariamente con la stessa chiave).
    \item La funzione può aggregare, filtrare, trasformare o combinare i valori.
  \end{itemize}
\end{definizione}

% ─────────────────────────────────────────────────────────────────────
\subsection{Fase Shuffle / Group-by-Key}

Dopo la fase map, tutte le coppie intermedie emesse passano alla fase
\textbf{shuffle/group-by-key}:
\begin{itemize}
  \item Tutte le coppie con la stessa chiave vengono \textbf{raggruppate}
    in una lista di valori.
  \item Formeranno l'input del reduce: una collezione di coppie
    $\kv{k}{\Lk}$.
  \item Molti runtime che implementano MapReduce \textbf{ordinano}
    anche le chiavi all'interno di ogni partizione.
\end{itemize}

\textbf{Perché il raggruppamento:} tra map task diversi (su nodi diversi)
molte coppie intermedie possono condividere la stessa chiave. Per la maggior
parte dei problemi, dobbiamo riunire tutti i valori che condividono una chiave
per calcolare un risultato globale.

\textbf{Come le chiavi vengono assegnate ai reduce task:}
\begin{itemize}
  \item Definito da una \emph{politica di partizionamento} (es.\ hashing
    sulla chiave), gestita dal runtime.
  \item Ogni reduce task elabora una chiave alla volta, ma può gestirne
    molte complessivamente.
  \item \textbf{Invariante chiave:} tutti i valori con la stessa chiave
    vengono inviati allo stesso reduce task.
\end{itemize}

% ─────────────────────────────────────────────────────────────────────
\subsection{Definizione rigorosa del job MapReduce}

\textbf{Input:} un multiset (o collezione) di coppie chiave-valore $I$.

\textbf{Scelte del progettista:}
\begin{enumerate}
  \item Definire le due funzioni utente: $\MAP$ e $\REDUCE$.
  \item Specificare i formati e i tipi di dati di input/output.
  \item Definire opzionalmente una politica di partizionamento personalizzata.
\end{enumerate}

\textbf{Responsabilità del runtime:}
\begin{itemize}
  \item Distribuire l'input tra i map task.
  \item Eseguire le funzioni map in parallelo.
  \item Fare shuffle e raggruppare i dati intermedi per chiave.
  \item Eseguire le funzioni reduce in parallelo.
  \item Gestire fallimenti e stragglers.
\end{itemize}

% ─────────────────────────────────────────────────────────────────────
\subsection{Sincronizzazione in MapReduce}

Un job MapReduce richiede punti di sincronizzazione nel cluster:
\begin{itemize}
  \item Elaborazione locale delle partizioni di input (map task).
  \item Ridistribuzione e raggruppamento dei dati intermedi per chiave (shuffle).
  \item Aggregazione per produrre gli output finali (reduce task).
\end{itemize}

Questo richiede coordinazione: le fasi avvengono in un ordine ben definito
(map $\to$ shuffle $\to$ reduce) e i dati intermedi devono essere spostati
e raggruppati.

Nel modello MapReduce \textbf{astraiamo via} i meccanismi: progettiamo il
modello dei dati $\kv{k}{v}$ e le due funzioni ($\MAP$/$\REDUCE$); il
runtime garantisce la corretta esecuzione sotto le assunzioni del modello.

% ─────────────────────────────────────────────────────────────────────
\subsection{Come pensare in MapReduce}

\textbf{Per risolvere un problema in MapReduce} seguiamo due passi di progetto:

\begin{enumerate}
  \item \textbf{Modellare i dati:}
    \begin{itemize}
      \item Rappresentare input e output come coppie chiave-valore.
      \item Scegliere le chiavi in modo che guidino il partizionamento e il
        raggruppamento (shuffle).
      \item Ogni dato riceve una chiave che guida storage, computazione
        e movimento.
    \end{itemize}
  \item \textbf{Progettare due funzioni:}
    \begin{itemize}
      \item \textbf{Map:} elabora gli elementi di input indipendentemente
        ed emette coppie chiave-valore intermedie.
      \item \textbf{Reduce:} aggrega ed elabora i valori che condividono
        la stessa chiave.
    \end{itemize}
\end{enumerate}

\textbf{Divisione delle responsabilità:}
\begin{itemize}
  \item Il \textbf{programmatore} specifica: funzioni Map e Reduce; file di
    input / sorgenti dati.
  \item Il \textbf{sistema} gestisce: lettura degli input come coppie
    chiave-valore; shuffle e sort verso i nodi di output; routing di tutte
    le coppie con una data chiave allo stesso reducer; scrittura dei risultati
    su file.
\end{itemize}

\begin{takeaway}
  Il programmatore si concentra su \emph{cosa} calcolare; il sistema gestisce
  \emph{come} distribuire e coordinare.
\end{takeaway}

% ─────────────────────────────────────────────────────────────────────
\newpage
\section{WordCount: Esempio Illustrativo}
% ─────────────────────────────────────────────────────────────────────

\subsection{Definizione del problema}

\begin{esempio}
  \textbf{Scenario:} l'input è una enorme collezione di documenti testuali;
  vogliamo contare le occorrenze delle parole.

  \textbf{Definizione formale:}
  \begin{itemize}
    \item \textbf{Input:} una collezione di documenti (ogni documento è una
      sequenza di parole/stringhe).
    \item \textbf{Output:} per ogni parola $w$, il numero totale di occorrenze
      di $w$ nell'intera collezione.
  \end{itemize}

  \textbf{Applicazione tipica:} analytics e machine learning su larga scala.
\end{esempio}

% ─────────────────────────────────────────────────────────────────────
\subsection{Passo 1: Modellazione dei dati}

\begin{itemize}
  \item \textbf{Input:} una collezione $D$ di documenti.
    \begin{itemize}
      \item Ogni documento è un \emph{elemento}.
      \item Un chunk è un insieme di elementi assegnato a un map task.
      \item Le chiavi di input possono essere ID documento; per WordCount
        non sono essenziali.
      \item Il valore è il contenuto del documento.
    \end{itemize}
  \item \textbf{Output:} per ogni parola, il totale delle occorrenze.
    \begin{itemize}
      \item Chiavi: parole (tipo \texttt{String}).
      \item Valori: conteggi di occorrenze (tipo \texttt{Integer}).
    \end{itemize}
  \item \textbf{Dati intermedi:} coppie della forma $\kv{w}{1}$.
    \begin{itemize}
      \item Emesse dalla funzione map per ogni occorrenza di parola.
      \item Raggruppate per chiave durante lo shuffle.
    \end{itemize}
\end{itemize}

% ─────────────────────────────────────────────────────────────────────
\subsection{Passo 2: Definizione della funzione Map}

La funzione map elabora un elemento di input (un documento) e produce
coppie chiave-valore intermedie.

\textbf{Algoritmo:}
\begin{enumerate}
  \item Legge un documento.
  \item Lo tokenizza in parole $w_1, w_2, \ldots, w_n$.
  \item Per ogni parola $w_i$, \emph{emette} la coppia $\kv{w_i}{1}$.
\end{enumerate}

\textbf{Significato:}
\begin{itemize}
  \item Chiave = la parola $w_i$.
  \item Valore = 1, ovvero ``un'occorrenza osservata''.
\end{itemize}

\textbf{Risultato dell'elaborazione di un chunk:} emissione di una
collezione di coppie
\[
  \kv{w_1}{1},\; \kv{w_2}{1},\; \ldots,\; \kv{w_n}{1}
\]
una coppia per ogni parola in ogni documento.

\begin{nota}
  Se la parola $m$ appare $k$ volte in un documento $D_i$, avremo $k$
  coppie $\kv{m}{1}$ nell'output emesso. I dati intermedi sono memorizzati
  localmente nel nodo worker, poi trasmessi durante lo shuffle.
\end{nota}

% ─────────────────────────────────────────────────────────────────────
\subsection{Passo 3: Definizione della funzione Reduce}

\textbf{Input del reduce:} coppie raggruppate $\kv{k}{[l_1, l_2, \ldots, l_n]}$
\begin{itemize}
  \item Per WordCount, ogni $l_i = 1$ (un'occorrenza).
  \item $n$ è il numero totale di volte in cui la parola $k$ è apparsa
    in tutti i documenti.
\end{itemize}

\textbf{Algoritmo:}
\begin{enumerate}
  \item Calcola $\texttt{COUNT}(k) = \sum_{i=1}^{n} l_i = n$.
  \item Emette (output) la coppia $\kv{k}{\texttt{COUNT}(k)}$.
\end{enumerate}

\textbf{Output finale:} ogni reducer emette coppie $\kv{\mathit{word}}{\mathit{count}}$
per le parole che ha elaborato. L'unione di tutti gli output dei reducer fornisce
il conteggio completo delle parole per l'intera corpus.

\begin{nota}
  Produrre un singolo file richiede un passo aggiuntivo (es.\ una fusione
  successiva o un job con un solo reducer).
\end{nota}

% ─────────────────────────────────────────────────────────────────────
\subsection{Pseudocodice MapReduce}

Il pseudocodice di un algoritmo MapReduce deve contenere:
\begin{itemize}
  \item \textbf{Input per map:} un elemento del chunk assegnato.
  \item \textbf{Input di reduce:} una coppia chiave-lista di valori,
    risultato del raggruppamento di tutti i valori-chiave generati dai map.
  \item \textbf{Corpo delle procedure:} definizione della computazione
    e degli output (tramite \texttt{emit}).
  \item \textbf{Extra (opzionale):} piccole variabili condivise/parametri,
    codice di inizializzazione o chiusura.
\end{itemize}

\begin{lstlisting}[language=Python, caption={WordCount -- Mapper}]
class Mapper:
    """Il mapper emette una coppia chiave-valore intermedia
       per ogni parola in un documento."""
    def Map(self, docid_a, doc_d):
        for term_t in doc_d:
            Emit(t, 1)
\end{lstlisting}

\begin{lstlisting}[language=Python, caption={WordCount -- Reducer}]
class Reducer:
    """Il reducer somma tutti i conteggi per ogni termine."""
    def Reduce(self, term_t, list_of_counts_Lt):
        total = 0
        for c in list_of_counts_Lt:
            total += c
        Emit(t, total)
\end{lstlisting}

% ─────────────────────────────────────────────────────────────────────
\newpage
\section{Principi di Progettazione}
% ─────────────────────────────────────────────────────────────────────

\subsection{Pro e contro di un modello rigido}

\textbf{Pro del modello rigido:} molte problematiche vengono gestite dal
sistema per design (partizionamento di default, sincronizzazione, ridondanza,
\ldots); semplicità, efficacia, affidabilità, facilità d'uso.

\textbf{Contro:} controllo diretto limitato su:
\begin{itemize}
  \item \emph{Dove} gira un mapper/reducer.
  \item \emph{Quando} un task inizia/finisce.
  \item \emph{Quali} coppie di input vengono elaborate da quale mapper.
  \item \emph{Quali} coppie intermedie vanno a quale reducer (solo:
    tutti i valori per una chiave vanno a un reducer).
\end{itemize}

\begin{importante}
  Le prestazioni sono controllabili solo \textbf{indirettamente} tramite
  i punti di personalizzazione consentiti.
\end{importante}

% ─────────────────────────────────────────────────────────────────────
\subsection{Fattori che influenzano le soluzioni MR}

L'efficacia e la correttezza delle soluzioni MR dipendono fortemente da:
\begin{itemize}
  \item \textbf{Tipo di problema:} è effettivamente parallelizzabile?
  \item \textbf{Qualità del design dell'algoritmo:} sfrutta efficacemente
    il parallelismo e supera un forte baseline sequenziale?
  \item \textbf{Il sistema sottostante:} hardware, rete e vincoli di storage;
    trade-off costo/prestazioni.
\end{itemize}

\begin{importante}
  \textbf{Requisito di correttezza:} gli algoritmi devono essere corretti
  indipendentemente da come i record di input sono divisi tra i mapper,
  e tutti i valori per la stessa chiave devono essere elaborati insieme
  da un reducer.
\end{importante}

% ─────────────────────────────────────────────────────────────────────
\subsection{Effetti di sistema sui cluster}

A differenza degli algoritmi sequenziali tradizionali, sui cluster:
\begin{itemize}
  \item Il tempo speso nella \textbf{comunicazione} tra nodi è
    non-trascurabile e deve essere modellato.
  \item \textbf{Trasmettere dati} spesso costa molto di più che elaborare
    la stessa quantità di dati localmente.
  \item Le risorse eterogenee delle macchine (es.\ memoria per nodo)
    devono essere considerate per evitare casi patologici.
\end{itemize}

% ─────────────────────────────────────────────────────────────────────
\subsection{Manopole di tuning in MapReduce}

\begin{tabular}{@{}lp{8.5cm}@{}}
  \toprule
  \textbf{Manopola} & \textbf{Descrizione} \\
  \midrule
  Parallelismo    & Regolare il numero di task e/o il numero di macchine. \\
  Shuffling       & Personalizzare il partizionamento e come i dati intermedi
                    sono distribuiti. \\
  Sincronizzazione & Decidere se tutti i task di una fase devono finire prima
                    che la successiva possa iniziare. \\
  Ridondanza      & Eseguire copie di backup dei task per tollerare
                    i fallimenti e mitigare gli stragglers. \\
  \bottomrule
\end{tabular}

% ─────────────────────────────────────────────────────────────────────
\subsection{Tecniche di personalizzazione}

Cosa può fare un programmatore/progettista di algoritmi per ottimizzare MR
senza rompere il modello:
\begin{itemize}
  \item \textbf{Usare chiavi/valori strutturati} (non solo tipi primitivi):
    cambia come i dati sono distribuiti ed elaborati.
  \item \textbf{Codice di inizializzazione/terminazione} nei task map/reduce:
    abilita stato condiviso locale al task tra chiamate di funzione.
  \item \textbf{Preservare variabili di stato} nei mapper/reducer tra
    input/chiavi: abilita riuso locale al nodo (attenzione: correttezza +
    determinismo).
  \item \textbf{Controllare l'ordinamento delle chiavi intermedie:}
    può aiutare la correttezza o le prestazioni in problemi specifici.
  \item \textbf{Personalizzare il partizionamento dello spazio delle chiavi:}
    controlla quale reducer vede quali chiavi (invariante: una chiave
    $\to$ un reducer).
  \item \textbf{Parametri ambientali:} es.\ numero di reducer, molteplicità
    dei task, ridondanza.
\end{itemize}

Quando una tecnica si applica a molte problematiche nella stessa categoria,
la chiamiamo \textbf{design pattern}.

% ─────────────────────────────────────────────────────────────────────
\newpage
\section{Framework di Esecuzione}
% ─────────────────────────────────────────────────────────────────────

\subsection{Cosa gestisce il sistema}

\begin{tabular}{@{}ll@{}}
  \toprule
  \textbf{Sistema (gestisce automaticamente)} & \textbf{Progettista (si concentra su)} \\
  \midrule
  Distribuzione e movimento dei dati (shuffle) & Modellazione dati (chiavi/valori) \\
  Scheduling e sincronizzazione dei task       & Logica algoritmica (Map/Reduce) \\
  Fault tolerance (ri-esecuzione su fallimenti) & Trade-off prestazionali \\
  Storage e replicazione su DFS                &  \\
  \bottomrule
\end{tabular}

Come progettisti ignoriamo: eterogeneità hardware, topologia di rete,
sistemi operativi e altri dettagli di basso livello.

% ─────────────────────────────────────────────────────────────────────
\subsection{Flusso computazionale (control plane)}

\begin{enumerate}
  \item Il programma utente avvia un \textbf{master controller}.
  \item Il programma utente avvia un insieme di processi \textbf{worker}.
\end{enumerate}

\begin{itemize}
  \item \textbf{Master:} coordina globalmente.
  \item \textbf{Worker:} esegue map task o reduce task.
\end{itemize}

% ─────────────────────────────────────────────────────────────────────
\subsection{Movimento dei dati}

\begin{enumerate}
  \item L'input risiede in un DFS (partizionato in split/blocchi).
  \item Lo scheduler piazza i task \textbf{vicino ai dati} (data locality).
  \item Gli output dei map sono memorizzati sui \textbf{dischi locali}
    dei worker.
  \item Durante lo shuffle, gli output dei map vengono trasferiti ai
    reducer in base al partitioner.
  \item L'output finale viene scritto di nuovo nel DFS (o in un sink esterno).
\end{enumerate}

\begin{nota}
  L'output di un job MR può diventare l'input di un altro
  (algoritmi multi-round).
\end{nota}

% ─────────────────────────────────────────────────────────────────────
\subsection{Responsabilità del Master Controller}

\begin{enumerate}
  \item \textbf{Avvio} dei map e reduce task.
  \item \textbf{Assegnazione} dei task ai worker.
  \item \textbf{Tracciamento} dello stato dei task, dei worker/nodi, del
    movimento dei dati (posizioni degli output intermedi, informazioni
    di partizionamento).
\end{enumerate}

% ─────────────────────────────────────────────────────────────────────
\subsection{Responsabilità dei Worker}

\begin{enumerate}
  \item Esecuzione dei task map/reduce (la computazione effettiva).
  \item Segnalazione del completamento del task al master.
  \item Produzione degli output intermedi/finali.
  \item Creazione di file locali quando necessario.
  \item Comunicazione al master della posizione e dimensione dei file prodotti.
\end{enumerate}

% ─────────────────────────────────────────────────────────────────────
\subsection{Algoritmo di scheduling dei task}

Ogni sistema MR include un algoritmo di scheduling per:
\begin{itemize}
  \item Assegnare i task ai worker in modo efficace.
  \item Gestire l'interazione worker-master.
\end{itemize}

\textbf{Stato di ogni task:}
\begin{itemize}
  \item \textbf{Idle:} in attesa di essere schedulato.
  \item \textbf{In-progress:} attualmente in esecuzione.
  \item \textbf{Completed:} segnala al master la posizione e la dimensione
    dei suoi file di output (gli output dei map sono tipicamente su dischi
    locali; gli output dei reduce sono scritti nel DFS).
\end{itemize}

Il master fa ping periodicamente ai worker per monitorare lo stato,
schedulare nuovi task e rilevare fallimenti.

% ─────────────────────────────────────────────────────────────────────
\subsection{Gestione dei fallimenti}

\textbf{In caso di fallimento del worker:}
\begin{itemize}
  \item I map task (completati o in-progress) sul worker fallito vengono
    resettati a idle e rischedulati, poiché anche gli output dei map
    completati diventano non disponibili se memorizzati localmente.
  \item I reduce task in-progress sul worker fallito vengono resettati
    a idle e rischedulati.
\end{itemize}

\textbf{Fallimento del Master:} l'intero job MR viene abortito e riavviato.

\begin{tabular}{@{}lll@{}}
  \toprule
  \textbf{Tipo di fallimento} & \textbf{Impatto} & \textbf{Soluzione} \\
  \midrule
  Master   & Worst case & Riavvio dell'intero job \\
  Map task & Solitamente minore & Riassegnazione e riavvio su altro worker \\
  Reduce task & Solitamente minore & Riassegnazione e riavvio \\
  \bottomrule
\end{tabular}

\textbf{Perché non mantenere lo stato parzialmente elaborato?}
\begin{itemize}
  \item Tracciare i progressi parziali aggiunge complessità implementativa
    significativa.
  \item Il checkpointing può introdurre overhead di CPU, memoria e I/O.
  \item Riavviare è spesso più economico e semplice del recovery fine-grained.
\end{itemize}

% ─────────────────────────────────────────────────────────────────────
\subsection{Granularità dei task}

Avere \textbf{più di un map task per nodo} è conveniente:
\begin{itemize}
  \item Se il Worker 1 completa due map task prima che il Worker 2
    completi un singolo task pesante (\emph{straggler}), i reduce task
    non possono iniziare finché tutti gli output map richiesti non sono
    disponibili.
  \item Avere molti task permette di bilanciare meglio il carico.
\end{itemize}

% ─────────────────────────────────────────────────────────────────────
\subsection{Backup task (esecuzione speculativa)}

\begin{definizione}[-- Backup Task]
  I backup task sono un metodo euristico e proattivo per ottenere output
  più velocemente:
  \begin{itemize}
    \item Si avviano copie ridondanti dei task a esecuzione lenta
      (per mitigare gli \emph{stragglers}).
    \item Una volta che la prima copia completa, le copie rimanenti
      vengono terminate.
  \end{itemize}
\end{definizione}

\textbf{Perché funziona:} i task possono bloccarsi a causa di dischi difettosi,
rallentamenti CPU transitori, rumore del sistema operativo, ecc. Si scambia
risorse computazionali extra per una latenza di coda (\emph{job completion time})
più bassa.

% ─────────────────────────────────────────────────────────────────────
\newpage
\section{Partizionamento e Meccanismi di Sistema}
% ─────────────────────────────────────────────────────────────────────

\subsection{Vista end-to-end: MapReduce su DFS}

\begin{itemize}
  \item \textbf{Input:} i dati vengono ingestiti nel DFS (il passo
    \texttt{copyFromLocal} è opzionale).
  \item \textbf{Fase Map:} task paralleli leggono gli input split
    (preferibilmente data-local) e generano coppie $\kv{k}{v}$ intermedie.
  \item \textbf{Fase Shuffle/Group:} le coppie intermedie vengono ridistribuite
    affinché i valori con la stessa chiave si incontrino.
  \item \textbf{Fase Reduce:} task paralleli elaborano i valori per chiave
    e producono coppie chiave-valore di output.
  \item \textbf{Output:} scritto di nuovo nel DFS; copiato fuori
    (\texttt{copyToLocal}) solo se necessario.
\end{itemize}

\begin{takeaway}
  Il DFS fornisce storage scalabile; MapReduce fornisce un modo strutturato
  per eseguire computazioni sui dati distribuiti.
\end{takeaway}

% ─────────────────────────────────────────────────────────────────────
\subsection{Partizionamento: panoramica}

\begin{importante}
  \textbf{Vincolo del modello:} una chiamata alla funzione reduce elabora
  esattamente una coppia raggruppata $\kv{k}{\Lk}$. Tutti i valori con la
  stessa chiave $k$ sono raccolti in $\Lk$.
\end{importante}

Tipicamente il numero di reduce task $r$ è configurabile dall'utente.

\textbf{Domanda sull'assegnazione delle chiavi:} come decidiamo quale reduce task
gestisce quali chiavi?
\begin{itemize}
  \item Se $r \geq$ numero di chiavi distinte: si potrebbe assegnare una
    chiave per reduce task (raramente fattibile).
  \item Altrimenti: le chiavi \textbf{devono essere partizionate} tra i
    reduce task.
\end{itemize}

% ─────────────────────────────────────────────────────────────────────
\subsection{Partizionamento tramite hashing}

Una politica di partizionamento comune usa una funzione hash sulle chiavi:
\[
  h: U \to \{0, 1, \ldots, r-1\}
\]
dove $U$ è l'universo delle chiavi possibili e $r$ è il numero di reduce task
(tipicamente $r \ll |U|$). La chiave $k$ è assegnata al reduce task n.~$h(k)$.

\begin{itemize}
  \item $h(k)$ è detto \textbf{numero di bucket} o \textbf{partizione}
    della chiave $k$.
  \item Tutte le chiavi con lo stesso numero di bucket vanno allo stesso
    reduce task.
\end{itemize}

\textbf{Proprietà di buone funzioni hash:}
\begin{itemize}
  \item \textbf{Deterministica:} la stessa chiave va sempre nello stesso bucket.
  \item \textbf{Uniforme:} le chiavi sono distribuite equamente tra i bucket
    (idealmente).
  \item \textbf{Veloce:} overhead computazionale minimo.
\end{itemize}

\textbf{Come i dati intermedi raggiungono i reduce task:}
\begin{enumerate}
  \item Ogni map task produce coppie intermedie e le scrive localmente.
  \item L'output del map è partizionato in $r$ partizioni usando $h(k)$.
  \item La partizione $i$ contiene coppie le cui chiavi soddisfano $h(k) = i$.
  \item La partizione $i$ verrà inviata al reduce task $i$.
  \item Molti runtime \textbf{ordinano} i dati per chiave all'interno di ogni
    partizione.
\end{enumerate}

% ─────────────────────────────────────────────────────────────────────
\subsection{Il problema dello skew}

\begin{esempio}
  Supponiamo $r = 4$ reduce task e chiavi intere in $[0, 2^{10}-1]$.
  Usando $h(k) = k \bmod 4$, ogni chiave è assegnata deterministicamente
  a un reduce task. Nel caso peggiore (tutte le $2^{10}$ chiavi presenti),
  ogni reducer può ricevere fino a $2^{10}/4 = 2^8 = 256$ chiavi distinte.

  \textbf{Caso patologico:} se tutte le chiavi sono multiple di 4, con
  $h(k) = k \bmod 4$ tutte vanno al reducer 0 $\Rightarrow$ un reducer
  fa (quasi) tutto il lavoro $\Rightarrow$ \textbf{parallelismo sprecato}.
\end{esempio}

\begin{definizione}[-- Skew]
  Lo \textbf{skew} si riferisce a uno squilibrio significativo tra i reduce task:
  \begin{itemize}
    \item Reducer diversi impiegano tempi diversi.
    \item Il tempo di completamento globale del job è spesso determinato
      dal \textbf{reducer più lento}.
  \end{itemize}
  Può derivare da:
  \begin{itemize}
    \item Una mappatura sbilanciata di chiavi su reducer (partizionamento sbagliato).
    \item Dimensioni delle liste di valori $|\Lk|$ molto diseguali
      (\emph{hot key}).
  \end{itemize}
\end{definizione}

% ─────────────────────────────────────────────────────────────────────
\subsection{Scelta della strategia di partizionamento}

\begin{itemize}
  \item Dipende dalla \textbf{distribuzione delle chiavi} nei dati.
  \item La maggior parte dei sistemi fornisce partitioner di default che
    funzionano bene per tipi di chiave comuni.
  \item Opzionalmente, l'utente può fornire un \textbf{partitioner
    personalizzato}.
\end{itemize}

\textbf{Trade-off nella scelta del numero di reduce task $r$:}
\begin{itemize}
  \item $r$ maggiore $\Rightarrow$ più parallelismo ma più overhead.
  \item $r$ minore $\Rightarrow$ meno overhead ma potenzialmente meno
    parallelismo.
  \item Usare più reduce task dei nodi può aiutare a mitigare lo skew:
    i reducer a lungo termine occupano alcuni nodi; i reducer più brevi
    possono girare su altri nodi, migliorando l'utilizzo.
\end{itemize}

% ─────────────────────────────────────────────────────────────────────
\newpage
\section{Esempio di Implementazione: Apache Hadoop}
% ─────────────────────────────────────────────────────────────────────

\subsection{Architettura classica di Hadoop}

A livello implementativo, la coordinazione è effettuata da un
\textbf{master service} che:
\begin{itemize}
  \item Schedula i task e monitora i progressi.
  \item Riavvia i task falliti.
  \item Coordina i trasferimenti dati durante lo shuffle.
\end{itemize}

Un sistema MapReduce funzionante è quindi uno \textbf{stack complesso di
componenti} che interagiscono.

% ─────────────────────────────────────────────────────────────────────
\subsection{Componenti principali dell'architettura Hadoop classica}

\begin{itemize}
  \item \textbf{Client Program:} invia un job (programma + configurazione)
    al sistema.
  \item \textbf{JobTracker} (master): coordina l'esecuzione.
    \begin{itemize}
      \item Divide l'input in input split.
      \item Assegna map task e reduce task ai worker.
      \item Monitora l'esecuzione e gestisce i fallimenti.
    \end{itemize}
  \item \textbf{TaskTracker} (worker daemon): ogni nodo worker esegue un
    TaskTracker.
    \begin{itemize}
      \item Esegue map task su split (preferibilmente locali).
      \item In seguito esegue reduce task sui dati intermedi raggruppati.
    \end{itemize}
  \item \textbf{Distributed File System (DFS):} memorizza input e output.
\end{itemize}

\begin{importante}
  Questo è un \textbf{esempio di implementazione} (una possibile architettura
  runtime). In questo corso, il modello MapReduce è trattato come
  un'astrazione che nasconde questi componenti di sistema.
\end{importante}

% ─────────────────────────────────────────────────────────────────────
\newpage
\section{Riepilogo del Capitolo}
% ─────────────────────────────────────────────────────────────────────

\begin{tabular}{@{}lp{9cm}@{}}
  \toprule
  \textbf{Argomento} & \textbf{Concetti chiave} \\
  \midrule
  Motivazione &
    Decomposizione (divide\&conquer, input/output, owner-computes, speculativa);
    MPCM; shared-nothing. \\[4pt]
  Modello dati &
    Coppie $\kv{k}{v}$; chunk / input split; elementi. \\[4pt]
  Funzione Map &
    $\MAP(\kv{k}{v}) \to \kv{k'}{v'}^*$; una chiamata per elemento;
    output intermedi scritti localmente. \\[4pt]
  Shuffle / Group &
    Raggruppa per chiave $\to \kv{k}{\Lk}$; gestito dal sistema;
    invariante: stessa chiave $\to$ stesso reducer. \\[4pt]
  Funzione Reduce &
    $\REDUCE(\kv{k}{\Lk}) \to \kv{k''}{v''}^*$; una chiamata per chiave;
    può aggregare/filtrare/trasformare. \\[4pt]
  WordCount &
    Map emette $\kv{w}{1}$; Reduce somma i conteggi. \\[4pt]
  Design principles &
    Manopole di tuning; tecniche di personalizzazione; design pattern. \\[4pt]
  Execution framework &
    Master/worker; scheduling; stati dei task; fault tolerance;
    backup task. \\[4pt]
  Partizionamento &
    Hash $h(k) \in \{0,\ldots,r-1\}$; skew; hot key; trade-off su $r$. \\[4pt]
  Hadoop &
    JobTracker (master), TaskTracker (worker), DFS; solo un esempio
    implementativo. \\
  \bottomrule
\end{tabular}

\vspace{1em}
\begin{importante}
  \textbf{Prossimi argomenti:}
  \begin{itemize}
    \item Metriche di performance: $R$, $M_L$, $M_A$, $M_C$.
    \item Combiner e aggregazione locale.
    \item Classi di complessità MR: $\mathcal{MRC}^i$.
    \item Pratica: PySpark, Spark UI.
  \end{itemize}
\end{importante}

% ─────────────────────────────────────────────────────────────────────
\section*{Riferimenti}
% ─────────────────────────────────────────────────────────────────────
\begin{itemize}
  \item Jeffrey Dean e Sanjay Ghemawat,
    \textit{MapReduce: Simplified Data Processing on Large Clusters}.
    \url{http://labs.google.com/papers/mapreduce.html}
  \item Sanjay Ghemawat, Howard Gobioff e Shun-Tak Leung,
    \textit{The Google File System}.
    \url{http://labs.google.com/papers/gfs.html}
\end{itemize}
